\chapter{Introduzione}
\label{ch:introduzione}

L'obiettivo del presente lavoro di tesi è quello di sviluppare un sistema in grado di generare automaticamente una mesh 3D partendo da un modello CAD 2D di una planimetria di una singola stanza o di un appartamento.

Per ogni modello CAD 2D che verrà letto in input dal sistema il risultato sarà una mesh 3D con vertici, indici, normali e coordinate di texture che possono essere immediatamente utilizzate per il rendering.

Il dominio del problema comprende i piani architettonici CAD 2D in formato DXF (Drawing Exchange Format). Per poter iniziare a lavorare sulla progettazione della pipeline di elaborazione è stato prima necessario raccogliere un dataset composto da modelli CAD 2D rappresentativi del dominio di applicazione, più di dieci modelli CAD sono stati raccolti da diversi siti web come \texttt{draftsperson.net}, \texttt{pincad.com} e \texttt{archweb.com}.
Questi modelli CAD erano inizialmente disponibili nel formato DWG e sono quindi stati convertiti in formato DXF e poi controllati con un software di editing come LibreCAD.

Mentre DWG e DXF sono standard de facto nel campo del CAD sono diversi tra loro per quanto riguarda l'approccio alla codifica: DWG è un formato binario proprietario, più efficiente nello spazio ma meno leggibile in quanto può essere letto da poche applicazioni, mentre DXF è un formato basato sul testo (in ASCII) sviluppato in modo tale da non presentare problemi di interoperabilità con i software di terzi. \cite{adobe_dwg_dxf} Tra i modelli disponibili ne sono stati selezionati alcuni come studi di caso.

\section{I problemi dei modelli CAD: eterogeneità e mancanza di standard}
\label{sec:problemi_modelli_CAD}

Il formato DXF non fornisce direttamente poligoni, bensì un insieme di entità geometriche disconnesse tra loro, quali \texttt{POLYLINE}, \texttt{LINE}, \texttt{ARC}, \texttt{BLOCK} ecc., organizzate in layer identificati da un nome. È qui che emerge la prima criticità: non esiste alcuno standard che imponga una convenzione univoca sui nomi dei layer. Uno stesso layer contenente i muri può essere chiamato, a seconda del modello, \texttt{WALLS}, \texttt{A-WALL}, \texttt{WALL EXTERNAL}, \texttt{WALL INTERNAL}, e così via.

Un secondo problema riguarda la rappresentazione geometrica degli elementi. Non esiste uno standard neppure per la rappresentazione consistente dei muri: alcuni modelli li descrivono tramite polilinee, altri tramite semplici segmenti. Isolando le sole primitive relative ai muri, inoltre, si osservano ulteriori problemi ricorrenti: segmenti che si intersecano impropriamente tra loro, e contorni di muri che risultano aperti anziché chiusi.

Anche le finestre presentano una notevole variabilità rappresentativa: alcuni modelli le descrivono come blocchi complessi composti da migliaia di segmenti, altri come semplici segmenti che collegano due muri, altri ancora come rettangoli definiti da quattro segmenti. Le porte, al contrario, tendono a essere rappresentate in modo più uniforme tra i vari modelli, tipicamente tramite archi; nel caso peggiore, tuttavia, anche una porta, in particolare quella d'ingresso, può essere rappresentata tramite polilinee o segmenti anziché tramite un arco.

Questa eterogeneità è la ragione principale per cui non è possibile progettare un algoritmo generale in grado di elaborare correttamente qualsiasi modello CAD. È necessaria, a monte dell'elaborazione automatica, una fase di preprocessing volta a rendere il modello quanto più consistente possibile (Capitolo \ref{ch:geometry_processing}).

\section{Assunzioni progettuali}
\label{sec:assunzioni_progettuali}

Per affrontare queste criticità vengono fatte delle assunzioni ragionevoli, che consentono di semplificare il problema:
\begin{enumerate}
	\item i layer relativi a pareti, porte e finestre devono avere nomi riconoscibili (ad esempio \texttt{A-WALL}, \texttt{A-DOOR}, \texttt{A-WINDOW}). L'utente è tenuto a rinominarli manualmente;
	\item i muri devono essere già rappresentati come facce chiuse (poligoni), dotate di proprio spessore; eventuali \emph{crack} o discontinuità nei contorni devono essere riparate manualmente prima dell'esecuzione mediante software di editing CAD (come LibreCAD);
	\item le porte e finestre sono considerati buchi, o discontinuità,  nelle pareti e verranno trattati come tali durante la fase di estrusione, con altezze differenziate
\end{enumerate}